\documentclass[aspectratio=169,10pt]{beamer}
% Shared Beamer setup for Fall 2026 course slides.
% Clean Cal Poly-inspired theme: no custom class files or uncommon packages.
\mode<presentation>{
  \usetheme{default}
  \usefonttheme{professionalfonts}
}

\definecolor{CalPolyGreen}{HTML}{154734}
\definecolor{CalPolyGold}{HTML}{BD8B13}
\definecolor{CalPolySage}{HTML}{7FA28F}
\definecolor{CalPolyMint}{HTML}{DDF1E7}
\definecolor{CalPolySand}{HTML}{A29F86}
\definecolor{CalPolyBeige}{HTML}{EFEEDE}
\definecolor{DeckInk}{HTML}{1F2933}
\definecolor{DeckMuted}{HTML}{5F6B7A}
\definecolor{DeckSoft}{HTML}{F7F7F5}

\setbeamersize{text margin left=0.55in,text margin right=0.55in}
\setbeamercolor{normal text}{fg=DeckInk,bg=white}
\setbeamercolor{structure}{fg=CalPolyGreen}
\setbeamercolor{title}{fg=white}
\setbeamercolor{frametitle}{fg=CalPolyGreen,bg=white}
\setbeamercolor{block title}{bg=CalPolySage,fg=white}
\setbeamercolor{block body}{bg=CalPolyMint,fg=black}
\setbeamercolor{alerted text}{fg=CalPolyGold}

\setbeamerfont{title}{size=\Huge,series=\bfseries}
\setbeamerfont{frametitle}{size=\Large,series=\bfseries}
\setbeamerfont{block title}{size=\normalsize,series=\bfseries}

\setbeamertemplate{navigation symbols}{}
\setbeamertemplate{itemize item}{\color{CalPolyGold}$\blacktriangleright$}
\setbeamertemplate{itemize subitem}{\color{CalPolyGreen}$\bullet$}
\setbeamertemplate{footline}{
  \hfill{\color{DeckMuted}\scriptsize\insertframenumber/\inserttotalframenumber}
  \hspace{0.18in}\vspace{0.12in}
}
\setbeamertemplate{frametitle}{
  \vspace{0.18in}
  \hspace{0.02in}\insertframetitle\par
  \vspace{0.08in}
  \noindent\makebox[\linewidth][l]{%
    {\color{CalPolyGold}\rule{0.55in}{2pt}}%
    {\color{CalPolyGreen}\rule{\dimexpr\linewidth-0.55in\relax}{2pt}}%
  }
}

\newcommand{\CourseNumber}{Course}
\newcommand{\CourseTitle}{Course Title}
\newcommand{\LectureNumber}{0}
\newcommand{\LectureTitle}{Lecture Title}
\newcommand{\LectureDate}{Date}
\newcommand{\CourseUnit}{Unit}

\newcommand{\coursenumber}[1]{\renewcommand{\CourseNumber}{#1}}
\newcommand{\coursetitle}[1]{\renewcommand{\CourseTitle}{#1}}
\newcommand{\lecturenumber}[1]{\renewcommand{\LectureNumber}{#1}}
\newcommand{\lecturetitle}[1]{\renewcommand{\LectureTitle}{#1}}
\newcommand{\lecturedate}[1]{\renewcommand{\LectureDate}{#1}}
\newcommand{\courseunit}[1]{\renewcommand{\CourseUnit}{#1}}

\author{Kirk Duran}
\institute{California Polytechnic State University, San Luis Obispo}
\date{\LectureDate}

\newcommand{\makedecktitle}{
  \begingroup
  \setbeamercolor{background canvas}{bg=CalPolyGreen}
  \begin{frame}[plain]
    \vfill
    \begin{center}
      {\usebeamerfont{title}\color{white}\LectureTitle\par}
      \vspace{0.18in}
      {\Large\color{white}Lecture \LectureNumber\par}
    \end{center}
    \vfill
    \noindent{\color{white}\rule{\linewidth}{1.2pt}}\par
    \vspace{0.08in}
    \noindent\colorbox{CalPolyGold}{\makebox[0.24\linewidth][c]{\color{white}\Large\CourseNumber}}
    \hspace{0.12in}{\color{white}\Large Kirk Duran}
  \end{frame}
  \endgroup
}

\newenvironment{keyidea}{\begin{block}{Key idea}}{\end{block}}
\newenvironment{activity}{\begin{block}{In class}}{\end{block}}
\newenvironment{cpdefinition}{\begin{block}{Definition}}{\end{block}}
\newenvironment{exampleblockcp}{
  \begingroup
  \setbeamercolor{block title}{bg=CalPolySand,fg=white}
  \setbeamercolor{block body}{bg=CalPolyBeige,fg=black}
  \begin{block}{Example}
}{
  \end{block}
  \endgroup
}

\newcommand{\imageplaceholder}[2][0.3in]{%
  \noindent\makebox[\linewidth][c]{%
    \fcolorbox{CalPolySage}{CalPolyBeige}{%
      \parbox[c][#1][c]{0.86\linewidth}{%
        \centering
        {\color{DeckMuted}\scriptsize Image placeholder: }%
        {\color{DeckInk}\scriptsize\ttfamily\detokenize{#2}}%
      }%
    }%
  }%
}

\newcommand{\sectionframe}[1]{
  \begingroup
  \setbeamercolor{background canvas}{bg=CalPolyGreen}
  \begin{frame}[plain]
    \vfill
    {\color{white}\LARGE\bfseries #1\par}
    \vspace{0.18in}
    {\color{CalPolyGold}\rule{0.22\paperwidth}{3pt}}
    {\color{white}\rule{0.62\paperwidth}{3pt}}
    \vfill
  \end{frame}
  \endgroup
}

\newcommand{\placeholdernote}{
  \begin{block}{Placeholder}
    Replace these bullets with the final lecture material, examples, and in-class activities.
  \end{block}
}

\usepackage[T1]{fontenc}
\usepackage{lmodern}
\usepackage{graphicx}
\usepackage{fancyvrb}
\usepackage{tabularx}
\graphicspath{{../assets/lecture-33/}}
\setbeamersize{text margin left=0.42in,text margin right=0.42in}
\newcommand{\lectureimage}[2][2.50in]{\includegraphics[width=\linewidth,height=#1,keepaspectratio]{#2}}
\coursenumber{CS 1001}
\coursetitle{Introduction to Computing}
\lecturenumber{33}
\lecturetitle{Input, Output, and String Processing}
\lecturedate{November 16, 2026}
\courseunit{13. Input, Files, and Robust Programs}
\begin{document}
\makedecktitle
\begin{frame}[fragile]{Today}
\begin{itemize}
\item Explain how input and output cross a program boundary.
\item Use print(), input(), and numeric conversion intentionally.
\item Process text with strip(), split(), and simple string methods.
\item Inspect command-line arguments using sys.argv.
\item Organize runnable scripts around main() and the name guard.
\end{itemize}
\end{frame}
\begin{frame}[fragile]{Connection to Earlier Material}
\begin{itemize}
\item Our functions already accept parameters and return results.
\item Strings are sequences: indexing, loops, and selection patterns still apply.
\item This lecture adds the program boundary: where values come from and where results go.
\end{itemize}
\end{frame}
\sectionframe{Part 1: The Console as a Program Boundary}
\begin{frame}[fragile]{Input, Processing, Output}
\begin{columns}[T,onlytextwidth]\begin{column}{0.49\textwidth}
\centering\lectureimage[2.45in]{console-io-path.png}
\end{column}\begin{column}{0.47\textwidth}
\begin{itemize}
\item Input supplies values to a running program.
\item Processing uses the Python logic we already know.
\item Output makes information visible to the person running it.
\end{itemize}
\end{column}\end{columns}
\end{frame}
\begin{frame}[fragile]{Three Concepts Worth Separating}
\begin{block}{Python / console example}\begin{Verbatim}[fontsize=\footnotesize]
def double(n: int) -> int:
    return n * 2

print(double(4))  # console displays 8
\end{Verbatim}
\end{block}
\begin{itemize}
\item A return value moves data from a function back to its caller.
\item Printing writes a human-readable message to the console.
\item Reading input obtains text from a user or the command that launched the program.
\end{itemize}
\end{frame}
\begin{frame}[fragile]{print() Produces Console Output}
\begin{block}{Python / console example}\begin{Verbatim}[fontsize=\footnotesize]
print("Hello")
print("CS", 1001)
print("red", "green", sep=" / ")
print("loading", end="...")
print("done")
\end{Verbatim}
\end{block}
\begin{block}{Expected output}\begin{Verbatim}[fontsize=\footnotesize]
Hello
CS 1001
red / green
loading...done
\end{Verbatim}
\end{block}
\end{frame}
\begin{frame}[fragile]{Formatting Readable Output}
\begin{block}{Python / console example}\begin{Verbatim}[fontsize=\footnotesize]
name = "Ava"
score = 87.456
print(f"{name} scored {score:.1f}%")
print(f"Score: {score:.2f}")
\end{Verbatim}
\end{block}
\begin{block}{Expected output}\begin{Verbatim}[fontsize=\footnotesize]
Ava scored 87.5%
Score: 87.46
\end{Verbatim}
\end{block}
\begin{itemize}
\item f-strings combine values and formatting into one string.
\item Formatting changes how a value is displayed, not the value itself.
\end{itemize}
\end{frame}
\begin{frame}[fragile]{Predict: Return Versus Print}
\begin{block}{Python / console example}\begin{Verbatim}[fontsize=\footnotesize]
def announce(value: int) -> None:
    print(value + 1)

result = announce(4)
print(result)
\end{Verbatim}
\end{block}
\begin{activity}What two lines appear on the console? What is stored in result?\end{activity}
\end{frame}
\begin{frame}[fragile]{Reveal: Console Output Is Not a Return Value}
\begin{block}{Python / console example}\begin{Verbatim}[fontsize=\footnotesize]
5
None
\end{Verbatim}
\end{block}
\begin{itemize}
\item announce prints 5; it does not return a useful value.
\item Without an explicit return, Python returns None.
\item A testable helper often returns a result; the outer program decides when to print it.
\end{itemize}
\end{frame}
\sectionframe{Part 2: Keyboard Input and String Parsing}
\begin{frame}[fragile]{input() Waits for a Response}
\begin{block}{Python / console example}\begin{Verbatim}[fontsize=\footnotesize]
name = input("Name: ")
print(f"Hello, {name}!")
\end{Verbatim}
\end{block}
\begin{itemize}
\item The prompt is displayed first.
\item Execution waits until the user enters a line and presses Enter.
\item The returned value is a string.
\end{itemize}
\end{frame}
\begin{frame}[fragile]{Even Numeric-Looking Input Is Text}
\begin{columns}[T,onlytextwidth]\begin{column}{0.49\textwidth}
\centering\lectureimage[2.45in]{input-string-conversion.png}
\end{column}\begin{column}{0.47\textwidth}
\begin{itemize}
\item input() always returns str.
\item float(...) or int(...) performs explicit numeric conversion.
\item The text "12.5" is not the number 12.5.
\end{itemize}
\end{column}\end{columns}
\end{frame}
\begin{frame}[fragile]{Converting Text to Numbers}
\begin{block}{Python / console example}\begin{Verbatim}[fontsize=\footnotesize]
raw = input("Enter a temperature: ")
temp = float(raw)
print(f"Twice that value is {temp * 2}")
\end{Verbatim}
\end{block}
\begin{itemize}
\item float("3.5") produces 3.5.
\item int("12") produces 12.
\item For now, assume the user enters valid numeric text; error handling comes later.
\end{itemize}
\end{frame}
\begin{frame}[fragile]{Common String Operations for Input}
\begin{block}{Python / console example}\begin{Verbatim}[fontsize=\footnotesize]
raw = "  apple BANANA  "
clean = raw.strip()
words = clean.split()
print(clean.lower())
print(words)
\end{Verbatim}
\end{block}
\begin{block}{Expected output}\begin{Verbatim}[fontsize=\footnotesize]
apple banana
['apple', 'BANANA']
\end{Verbatim}
\end{block}
\begin{itemize}
\item strip() removes leading and trailing whitespace.
\item split() separates whitespace-delimited words.
\item lower() creates a normalized string.
\end{itemize}
\end{frame}
\begin{frame}[fragile]{Visual Trace: split()}
\begin{columns}[T,onlytextwidth]\begin{column}{0.49\textwidth}
\centering\lectureimage[2.45in]{split-tokens.png}
\end{column}\begin{column}{0.47\textwidth}
\begin{itemize}
\item Splitting creates a list of strings.
\item Repeated spaces become separators, not empty tokens, with split().
\item Numeric conversion is still a separate step.
\end{itemize}
\end{column}\end{columns}
\end{frame}
\begin{frame}[fragile]{Decompose a Pair Parser}
\begin{block}{Python / console example}\begin{Verbatim}[fontsize=\footnotesize]
def sum_pair(line: str) -> float:
    pieces = line.split()
    left = float(pieces[0])
    right = float(pieces[1])
    return left + right

print(sum_pair(" 3.5   -2.0 "))
\end{Verbatim}
\end{block}
\begin{block}{Expected output}\begin{Verbatim}[fontsize=\footnotesize]
1.5
\end{Verbatim}
\end{block}
\end{frame}
\begin{frame}[fragile]{Predict: Token Types}
\begin{block}{Python / console example}\begin{Verbatim}[fontsize=\footnotesize]
pieces = "12 8".split()
print(pieces[0] + pieces[1])
print(int(pieces[0]) + int(pieces[1]))
\end{Verbatim}
\end{block}
\begin{activity}What gets printed on each line? Why are the outputs different?\end{activity}
\end{frame}
\begin{frame}[fragile]{Reveal: Concatenation Versus Arithmetic}
\begin{block}{Expected output}\begin{Verbatim}[fontsize=\footnotesize]
128
20
\end{Verbatim}
\end{block}
\begin{itemize}
\item Both split() elements are str values.
\item The + operator concatenates strings.
\item Conversion is necessary before numeric addition.
\end{itemize}
\end{frame}
\sectionframe{Part 3: Interactive String Processing}
\begin{frame}[fragile]{Build a Useful Message from Input}
\begin{block}{Python / console example}\begin{Verbatim}[fontsize=\footnotesize]
first = input("First name: ").strip()
last = input("Last name: ").strip()
full_name = f"{first} {last}"
print(full_name.title())
\end{Verbatim}
\end{block}
\begin{itemize}
\item Read text, normalize it, compute a result, then display.
\item Strings are immutable; operations create new values.
\end{itemize}
\end{frame}
\begin{frame}[fragile]{An Input Loop with a Sentinel}
\begin{block}{Python / console example}\begin{Verbatim}[fontsize=\footnotesize]
def keyboard_sum() -> float:
    total = 0.0
    while True:
        text = input("Number (done to stop): ")
        if text == "done":
            return total
        total += float(text)
\end{Verbatim}
\end{block}
\begin{itemize}
\item The sentinel is a special value that ends the loop.
\item Check for the sentinel before numeric conversion.
\item The accumulator summarizes the numbers seen so far.
\end{itemize}
\end{frame}
\begin{frame}{Visual Trace: Sentinel Input}
\begin{center}
\lectureimage[1.85in]{sentinel-input-loop.png}
\end{center}
\begin{keyidea}Check the sentinel before converting; accumulate only numeric input.\end{keyidea}
\end{frame}
\begin{frame}[fragile]{Activity: Follow the Keyboard Session}
\begin{block}{Python / console example}\begin{Verbatim}[fontsize=\footnotesize]
Number (done to stop): 2.5
Number (done to stop): -1
Number (done to stop): 4
Number (done to stop): done
\end{Verbatim}
\end{block}
\begin{activity}What does keyboard\_sum() return? When does the last input get converted?\end{activity}
\end{frame}
\begin{frame}[fragile]{Answer: Accumulation Ends at done}
\begin{itemize}
\item The running totals are 2.5, 1.5, and 5.5.
\item The string "done" is checked, not converted.
\item The function returns 5.5.
\end{itemize}
\end{frame}
\sectionframe{Part 4: Command-Line Arguments}
\begin{frame}[fragile]{Arguments Are Supplied at Program Launch}
\begin{columns}[T,onlytextwidth]\begin{column}{0.49\textwidth}
\centering\lectureimage[2.45in]{argv-index-map.png}
\end{column}\begin{column}{0.47\textwidth}
\begin{itemize}
\item The command line can supply values before the program starts.
\item The interpreter launches the script with the provided arguments.
\item Unlike input(), these values do not require prompting during execution.
\end{itemize}
\end{column}\end{columns}
\end{frame}
\begin{frame}[fragile]{Inspect sys.argv}
\begin{block}{Python / console example}\begin{Verbatim}[fontsize=\footnotesize]
import sys

print(sys.argv)

# python total.py 3.5 2 -1
# ["total.py", "3.5", "2", "-1"]
\end{Verbatim}
\end{block}
\begin{itemize}
\item sys.argv is a list of strings.
\item Index 0 identifies the script, possibly including its path.
\item Indexes 1 onward hold the supplied arguments.
\end{itemize}
\end{frame}
\begin{frame}[fragile]{Numbers in Arguments Are Still Strings}
\begin{block}{Python / console example}\begin{Verbatim}[fontsize=\footnotesize]
import sys

# python echo.py 6 7
a = sys.argv[1]
b = sys.argv[2]
print(a + b)      # 67
print(int(a) + int(b))  # 13
\end{Verbatim}
\end{block}
\end{frame}
\begin{frame}[fragile]{Avoid Missing-Argument Index Errors}
\begin{block}{Python / console example}\begin{Verbatim}[fontsize=\footnotesize]
import sys

if len(sys.argv) < 2:
    print("Please provide a name")
else:
    print(f"Hello, {sys.argv[1]}!")
\end{Verbatim}
\end{block}
\begin{itemize}
\item Check length before indexing optional positions.
\item len(sys.argv) includes the script name.
\item Argument validation is a separate concern from processing.
\end{itemize}
\end{frame}
\begin{frame}[fragile]{Loop Over User Arguments}
\begin{block}{Python / console example}\begin{Verbatim}[fontsize=\footnotesize]
import sys

# python total.py 3.5 2 -1
total = 0.0
for text in sys.argv[1:]:
    total += float(text)
print(total)   # 4.5
\end{Verbatim}
\end{block}
\begin{itemize}
\item The slice [1:] skips the script name.
\item Each argument is converted before arithmetic.
\item Assume valid numeric arguments today.
\end{itemize}
\end{frame}
\begin{frame}[fragile]{Quotation Marks Group Arguments}
\begin{block}{Python / console example}\begin{Verbatim}[fontsize=\footnotesize]
$ python greet.py "Ada Lovelace"

# sys.argv[1] -> "Ada Lovelace"
\end{Verbatim}
\end{block}
\begin{itemize}
\item Without quotes, Ada and Lovelace are normally separate shell arguments.
\item The shell groups quoted text before Python receives sys.argv.
\item The shell syntax may look slightly different across operating systems.
\end{itemize}
\end{frame}
\begin{frame}[fragile]{Keyboard Input and CLI Arguments}
\begin{itemize}
\item input(): ask for a value after execution starts.
\item sys.argv: values are supplied when execution starts.
\item Both deliver strings and usually require parsing.
\item Choose the interface based on whether the program is interactive.
\end{itemize}
\end{frame}
\sectionframe{Part 5: main() and Organized Programs}
\begin{frame}[fragile]{A main() Function Expresses the Entry Point}
\begin{block}{Python / console example}\begin{Verbatim}[fontsize=\footnotesize]
def main() -> None:
    name = input("Name: ")
    print(f"Hello, {name}!")

main()
\end{Verbatim}
\end{block}
\begin{itemize}
\item main is a normal function name, not a Python keyword.
\item Organizing orchestration in main makes the flow easier to follow.
\item Calling main explicitly starts that part of the program.
\end{itemize}
\end{frame}
\begin{frame}[fragile]{Why the Name Guard Exists}
\begin{columns}[T,onlytextwidth]\begin{column}{0.49\textwidth}
\centering\lectureimage[2.45in]{main-entry-gate.png}
\end{column}\begin{column}{0.47\textwidth}
\begin{itemize}
\item When the script runs directly, \_\_name\_\_ has the value "\_\_main\_\_".
\item When imported, \_\_name\_\_ identifies the module instead.
\item The guard prevents the interactive program from automatically running on import.
\end{itemize}
\end{column}\end{columns}
\end{frame}
\begin{frame}[fragile]{The Standard Script Pattern}
\begin{block}{Python / console example}\begin{Verbatim}[fontsize=\footnotesize]
def main() -> None:
    print("Program started")

if __name__ == "__main__":
    main()
\end{Verbatim}
\end{block}
\begin{itemize}
\item Define functions first; call the entry point at the bottom.
\item Tests and other modules can import helpers without triggering the script.
\item The guard checks module execution context, not user input.
\end{itemize}
\end{frame}
\begin{frame}[fragile]{Predict: Running Versus Importing}
\begin{block}{Python / console example}\begin{Verbatim}[fontsize=\footnotesize]
# greetings.py
def greet(name: str) -> str:
    return f"Hi, {name}"

def main() -> None:
    print(greet("Ava"))

if __name__ == "__main__":
    main()
\end{Verbatim}
\end{block}
\begin{activity}What appears for python greetings.py? What appears for import greetings?\end{activity}
\end{frame}
\begin{frame}[fragile]{Reveal: The Guard Controls Side Effects}
\begin{itemize}
\item Running python greetings.py prints Hi, Ava.
\item Importing greetings defines greet and main without calling main.
\item The importing program can call greetings.greet("Ava") whenever it chooses.
\end{itemize}
\end{frame}
\begin{frame}[fragile]{Separate Calculation from Console I/O}
\begin{columns}[T,onlytextwidth]\begin{column}{0.49\textwidth}
\centering\lectureimage[2.45in]{program-architecture.png}
\end{column}\begin{column}{0.47\textwidth}
\begin{itemize}
\item A pure helper accepts values and returns a result.
\item main obtains inputs, invokes helpers, and displays results.
\item Separating responsibilities simplifies reasoning and unit tests.
\end{itemize}
\end{column}\end{columns}
\end{frame}
\begin{frame}[fragile]{Integrated Example: CLI Number Sum}
\begin{block}{Python / console example}\begin{Verbatim}[fontsize=\scriptsize]
import sys

def sum_numbers(text_values: list[str]) -> float:
    total = 0.0
    for text in text_values:
        total += float(text)
    return total

def main() -> None:
    total = sum_numbers(sys.argv[1:])
    print(f"Total: {total}")

if __name__ == "__main__":
    main()
\end{Verbatim}
\end{block}
\begin{itemize}
\item Run: python total.py 3.5 2 -1
\item Expected console output: Total: 4.5
\end{itemize}
\end{frame}
\begin{frame}[fragile]{Think-Pair-Share: Find the Bug}
\begin{block}{Python / console example}\begin{Verbatim}[fontsize=\footnotesize]
import sys

def main():
    amount = float(sys.argv[0])
    print(f"Amount: {amount}")

if __name__ == "__main__":
    main()
\end{Verbatim}
\end{block}
\begin{activity}What is wrong? What two changes would make this safer for a single numeric argument?\end{activity}
\end{frame}
\begin{frame}[fragile]{Reveal: Select and Check the Correct Argument}
\begin{block}{Python / console example}\begin{Verbatim}[fontsize=\footnotesize]
import sys

if len(sys.argv) < 2:
    print("Missing numeric argument")
else:
    amount = float(sys.argv[1])
    print(f"Amount: {amount}")
\end{Verbatim}
\end{block}
\begin{itemize}
\item sys.argv[0] is the script name, not the first user value.
\item Check that sys.argv[1] exists before indexing.
\item Full handling of malformed numeric text comes in a later lesson.
\end{itemize}
\end{frame}
\begin{frame}[fragile]{Common Wrong Models}
\begin{itemize}
\item “input() gives me a number” — it always returns a string.
\item “print() is the same as return” — one displays, one sends a value to a caller.
\item “sys.argv[0] is my first argument” — it identifies the script.
\item “main runs automatically” — it runs only when called.
\item “split() converts tokens” — it separates text; conversion is another step.
\end{itemize}
\end{frame}
\begin{frame}[fragile]{Takeaways}
\begin{itemize}
\item Programs exchange text with users through console input and output.
\item Parse strings explicitly before numeric computation.
\item Use strip(), split(), and loops to process user-entered text.
\item Use sys.argv[1:] for command-line values.
\item Put program orchestration in main() and protect it with the name guard.
\item Next lecture: input and output beyond a live console session.
\end{itemize}
\end{frame}
\end{document}
